% GENERATED FILE. Edit canonical lessons, then run npm run generate:books.
% canonical-source-hash: fnv1a64:074fd48fdda7a594
% canonical-lessons: JA-C462-shoubu, JA-C462-hikiwake, JA-C462-kyousou, JA-C462-jun-i, JA-C462-hansoku

\chapter{Win, Lose or Tie}
\label{ch:ja-win-lose-or-tie}
\hlchaptermodality{462}

\begin{chapteropening}
\textbf{By the end of this chapter:} \emph{I can talk about contests, rankings and fouls.}

Complete the last lesson of chapter 462: I can talk about contests, rankings and fouls.
\end{chapteropening}

\section[\texorpdfstring{\ja{しょうぶ} (shōbu)}{しょうぶ (shōbu)}]{\ja{しょうぶ} (shōbu) — a contest; a match}

\label{lesson:JA-C462-shoubu}



\begin{quote}
Before the new one: say the Japanese for gymnastics, then the Japanese for a team manager.
\end{quote}

\subsection*{You'll want to know: \ja{しょうぶ}}
\textbf{\ja{しょうぶ}} — \emph{shōbu} — \textquotedblleft{}a contest; a match\textquotedblright{}.

\begin{quote}
\textbf{\ja{あには} \ja{おとうとと} \ja{しょうぶを} \ja{しました。}} — \emph{ani wa otōto to shōbu o shimashita} — \textquotedblleft{}My elder brother had a contest with my younger brother.\textquotedblright{}
\end{quote}

\begin{cousinweb}
\textbf{\ja{しょうぶ}} (\emph{shōbu}) is Sino-Japanese (kango): two Chinese characters read with their on'yomi, readings borrowed from Chinese. They are \emph{shō}, \textquotedblleft{}victory\textquotedblright{} (Middle Chinese \emph{syingH}) and \emph{bu}, \textquotedblleft{}defeat\textquotedblright{} (\emph{bjuwX}) (KANJIDIC2; Middle Chinese in Baxter's transcription).
\end{cousinweb}

\begin{grammarlens}[title={What you've built}]
One more word for describing what you lived and felt.
\end{grammarlens}

\subsection*{Your turn}
\begin{itemize}
\raggedright
  \item \emph{Say it:} \emph{shōbu}
  \item \emph{Say it:} \emph{shōbu}, once more
  \item \emph{Say it:} \emph{ani wa otōto to shōbu o shimashita} — My elder brother had a contest with my younger brother.
  \item \emph{Recall:} say the Japanese for gymnastics, then the Japanese for a team manager, then say \emph{shōbu} again
\end{itemize}

\begin{tcolorbox}[breakable,colback=teal!4,colframe=teal!35!black,title={Before you move on}]
What does \ja{しょうぶ} mean? (\textquotedblleft{}A contest; a match\textquotedblright{}.) Say it once more.
\end{tcolorbox}
\section[\texorpdfstring{\ja{ひきわけ} (hikiwake)}{ひきわけ (hikiwake)}]{\ja{ひきわけ} (hikiwake) — a tie (in a game)}

\label{lesson:JA-C462-hikiwake}



\begin{quote}
Before the new one: say the Japanese for a team manager, then the Japanese for a contest.
\end{quote}

\subsection*{You'll want to know: \ja{ひきわけ}}
\textbf{\ja{ひきわけ}} — \emph{hikiwake} — \textquotedblleft{}a tie (in a game)\textquotedblright{}.

\begin{quote}
\textbf{\ja{しあいは} \ja{ひきわけでした。}} — \emph{shiai wa hikiwake deshita} — \textquotedblleft{}The match was a tie.\textquotedblright{}
\end{quote}

\begin{cousinweb}
\textbf{\ja{ひきわけ}} (\emph{hikiwake}) is a native Japanese compound (wago), inherited rather than borrowed from Chinese: \emph{hiki}, \textquotedblleft{}pulling (hiku)\textquotedblright{} and \emph{wake}, \textquotedblleft{}dividing (wakeru)\textquotedblright{} (JMdict, KANJIDIC2).
\end{cousinweb}

\begin{grammarlens}[title={What you've built}]
One more word for describing what you lived and felt.
\end{grammarlens}

\subsection*{Your turn}
\begin{itemize}
\raggedright
  \item \emph{Say it:} \emph{hikiwake}
  \item \emph{Say it:} \emph{hikiwake}, once more
  \item \emph{Say it:} \emph{shiai wa hikiwake deshita} — The match was a tie.
  \item \emph{Recall:} say the Japanese for a team manager, then the Japanese for a contest, then say \emph{hikiwake} again
\end{itemize}

\begin{tcolorbox}[breakable,colback=teal!4,colframe=teal!35!black,title={Before you move on}]
What does \ja{ひきわけ} mean? (\textquotedblleft{}A tie (in a game)\textquotedblright{}.) Say it once more.
\end{tcolorbox}
\section[\texorpdfstring{\ja{きょうそう} (kyōsō)}{きょうそう (kyōsō)}]{\ja{きょうそう} (kyōsō) — competition; rivalry}

\label{lesson:JA-C462-kyousou}



\begin{quote}
Before the new one: say the Japanese for a contest, then the Japanese for a tie.
\end{quote}

\subsection*{You'll want to know: \ja{きょうそう}}
\textbf{\ja{きょうそう}} — \emph{kyōsō} — \textquotedblleft{}competition; rivalry\textquotedblright{}.

\begin{quote}
\textbf{\ja{みせの} \ja{きょうそうが} \ja{きびしいです。}} — \emph{mise no kyōsō ga kibishii desu} — \textquotedblleft{}Competition between shops is fierce.\textquotedblright{}
\end{quote}

\begin{cousinweb}
\textbf{\ja{きょうそう}} (\emph{kyōsō}) is Sino-Japanese (kango): two Chinese characters read with their on'yomi, readings borrowed from Chinese. They are \emph{kyō}, \textquotedblleft{}compete with\textquotedblright{} (Middle Chinese \emph{gjaengH}) and \emph{sō}, \textquotedblleft{}contend\textquotedblright{} (\emph{tsreang}) (KANJIDIC2; Middle Chinese in Baxter's transcription).
\end{cousinweb}

\begin{grammarlens}[title={What you've built}]
One more word for describing what you lived and felt.
\end{grammarlens}

\subsection*{Your turn}
\begin{itemize}
\raggedright
  \item \emph{Say it:} \emph{kyōsō}
  \item \emph{Say it:} \emph{kyōsō}, once more
  \item \emph{Say it:} \emph{mise no kyōsō ga kibishii desu} — Competition between shops is fierce.
  \item \emph{Recall:} say the Japanese for a contest, then the Japanese for a tie, then say \emph{kyōsō} again
\end{itemize}

\begin{tcolorbox}[breakable,colback=teal!4,colframe=teal!35!black,title={Before you move on}]
What does \ja{きょうそう} mean? (\textquotedblleft{}Competition; rivalry\textquotedblright{}.) Say it once more.
\end{tcolorbox}
\section[\texorpdfstring{\ja{じゅんい} (jun'i)}{じゅんい (jun'i)}]{\ja{じゅんい} (jun'i) — a ranking; a placing}

\label{lesson:JA-C462-jun-i}



\begin{quote}
Before the new one: say the Japanese for a tie, then the Japanese for competition.
\end{quote}

\subsection*{You'll want to know: \ja{じゅんい}}
\textbf{\ja{じゅんい}} — \emph{jun'i} — \textquotedblleft{}a ranking; a placing\textquotedblright{}.

\begin{quote}
\textbf{\ja{しあいの} \ja{あとで、じゅんいが} \ja{きまりました。}} — \emph{shiai no ato de, jun'i ga kimarimashita} — \textquotedblleft{}After the match the rankings were decided.\textquotedblright{}
\end{quote}

\begin{cousinweb}
\textbf{\ja{じゅんい}} (\emph{jun'i}) is Sino-Japanese (kango): two Chinese characters read with their on'yomi, readings borrowed from Chinese. They are \emph{jun}, \textquotedblleft{}order\textquotedblright{} (Middle Chinese \emph{zywinH}) and \emph{i}, \textquotedblleft{}rank\textquotedblright{} (\emph{hwijH}) (KANJIDIC2; Middle Chinese in Baxter's transcription).
\end{cousinweb}

\begin{grammarlens}[title={What you've built}]
One more word for describing what you lived and felt.
\end{grammarlens}

\subsection*{Your turn}
\begin{itemize}
\raggedright
  \item \emph{Say it:} \emph{jun'i}
  \item \emph{Say it:} \emph{jun'i}, once more
  \item \emph{Say it:} \emph{shiai no ato de, jun'i ga kimarimashita} — After the match the rankings were decided.
  \item \emph{Recall:} say the Japanese for a tie, then the Japanese for competition, then say \emph{jun'i} again
\end{itemize}

\begin{tcolorbox}[breakable,colback=teal!4,colframe=teal!35!black,title={Before you move on}]
What does \ja{じゅんい} mean? (\textquotedblleft{}A ranking; a placing\textquotedblright{}.) Say it once more.
\end{tcolorbox}
\section[\texorpdfstring{\ja{はんそく} (hansoku)}{はんそく (hansoku)}]{\ja{はんそく} (hansoku) — a foul (in sport)}

\label{lesson:JA-C462-hansoku}



\begin{quote}
Before the new one: say the Japanese for competition, then the Japanese for a ranking.
\end{quote}

\subsection*{You'll want to know: \ja{はんそく}}
\textbf{\ja{はんそく}} — \emph{hansoku} — \textquotedblleft{}a foul (in sport)\textquotedblright{}.

\begin{quote}
\textbf{\ja{せんしゅが} \ja{はんそくを} \ja{しました。}} — \emph{senshu ga hansoku o shimashita} — \textquotedblleft{}A player committed a foul.\textquotedblright{}
\end{quote}

\begin{cousinweb}
\textbf{\ja{はんそく}} (\emph{hansoku}) is Sino-Japanese (kango): two Chinese characters read with their on'yomi, readings borrowed from Chinese. They are \emph{han}, \textquotedblleft{}anti-\textquotedblright{} (Middle Chinese \emph{pjonX}) and \emph{soku}, \textquotedblleft{}rule\textquotedblright{} (\emph{tsok}) (KANJIDIC2; Middle Chinese in Baxter's transcription).
\end{cousinweb}

\begin{grammarlens}[title={What you've built}]
One more word for describing what you lived and felt.
\end{grammarlens}

\subsection*{Your turn}
\begin{itemize}
\raggedright
  \item \emph{Say it:} \emph{hansoku}
  \item \emph{Say it:} \emph{hansoku}, once more
  \item \emph{Say it:} \emph{senshu ga hansoku o shimashita} — A player committed a foul.
  \item \emph{Recall:} say the Japanese for competition, then the Japanese for a ranking, then say \emph{hansoku} again
\end{itemize}

\begin{tcolorbox}[breakable,colback=teal!4,colframe=teal!35!black,title={Before you move on}]
What does \ja{はんそく} mean? (\textquotedblleft{}A foul (in sport)\textquotedblright{}.) Say it once more.
\end{tcolorbox}
